% abstract.tex — c01｜原文 p1｜论文标题页上的 Abstract（摘要）
% 0 条编号公式。
% [原文 p.1]
\noindent{\heiti\Large 摘要}\par
\addcontentsline{toc}{section}{摘要}
\vspace{0.4em}

在这些讲义中，我们给出超共形指标（superconformal index）的一个教学式导引。本稿是 2020 年冬季学校「YRISW 2020」上所作报告的文字版，将发表于 JPhysA 的一个专辑。讲义处于基础层次，面向有兴趣研究超共形指标的初年级研究生。

% [Tutor Note: 「教学式导引」对应 pedagogical introduction；作者有意面向初学者，正文保留大量口语与第一人称，译文不「润色」掉这种语气。]
% [Tutor Note: 术语 superconformal index 在 glossary 中冻结为「超共形指标」，不要写作「超共形指数」，以免与指数函数混淆。]
% [Tutor Note: YRISW 2020 是这次冬季学校的简称，原文未给全称；本译注本不擅自补全。]
% [Tutor Note: JPhysA 即 Journal of Physics A: Mathematical and Theoretical，原文说本稿将刊于其专辑。]
% [Tutor Note: 底本为 arXiv:2006.13630v1（2020 年 6 月 24 日，32 页）；本译本页标记用 PDF 页序，故正文第 1 页是 p.2。]
% [Tutor Note: 摘要只出现在 JHEP 标题页；正文 §0 的「讲义提纲」是另一段文字，二者不要混为一谈。]
% [Tutor Note: 原文 "interested in working with the superconformal index" 强调动手计算，译文用「研究」略失分寸，可理解为「以上手使用为目标」。]
% [Tutor Note: 摘要与 §0 都不带编号；§0 在原文的节计数器中算第 0 节，main.tex 用 \setcounter{section}{-1} 实现。]
